% Part: applied-modal-logics
% Chapter: temporal-logic
% Section: possible-histories

\documentclass[../../../include/open-logic-section]{subfiles}

\begin{document}

\olfileid{aml}{tl}{poss}

\olsection{ممکنې تاريخچې}

د زماني منطق هغه اړيکيز مدلونه چې تر اوسه مو کارولي، ډېر انعطاف لري؛
ځکه موږ پر لاسرسي اړيکه هېڅ قيد نۀ ږدو۔ نو زماني مدلونه په تېر او
راتلونکي دواړو لوريو کښې څانګې لرلے شي۔ خو ښايي وغواړو د څانګې يو
لا «مودالي» تصور وڅېړو چې پکښې د پېښو لړۍ د ممکنو تاريخچو په توګه
ګڼو۔ د دې دپاره د ژبې بدلول لازمي نۀ دي، که څۀ هم خپل «عادي» مودالي
عاملونه $\Box$ او~$\Diamond$ هم ورزياتولے شو؛ همدارنګه د پوهېدنې
د لاسرسي اړيکې ورزياتولے شو، څو د وخت په تېرېدو سره د عاملانو د
پوهې بدلون وښيو۔

\begin{defn}
  د زماني ژبې يو \emph{د ممکنو تاريخچو مدل} درې‌ګونې
  $\mModel{M} = \tuple{T, C, V}$ ده، چې په دې کښې
  \begin{enumerate}
    \item $T$ يو نا تش سټ دے، چې د وخت د حالتونو په توګه تفسيرېږي۔
    \item $C$ د يوه نظام د محاسبوي لارو، يا \emph{ممکنو
      تاريخچو}، سټ دے۔ په بله وينا، $C$ د حالتونو $s_1$، $s_2$،
      ~$s_3$، \dots د لړيو~$\sigma$ سټ دے، چې هر $s_i \in T$ وي۔
    \item $V$ داسې تابع ده چې هر !!{propositional
      variable}~$p$ ته د وخت د ټکو يو سټ~$V(p)$ ټاکي۔
  \end{enumerate}
  د اسانتيا دپاره عموماً دا هم فرضوو چې که يوه تاريخچه په~$C$ کښې
  وي، نو د هغې ټولې وروستۍ برخې هم پکښې وي۔ د بېلګې په توګه،
  که $s_1$، $s_2$،~$s_3$ په~$C$ کښې يوه لړۍ وي، نو $s_2$، $s_3$
  او~$s_3$ هم پکښې وي۔ همدارنګه، که دوه حالتونه $s_i$ او~$s_j$
  په لړۍ~$\sigma$ کښې راشي، نو د $i < j$ په حالت کښې وايو چې
  $s_i \prec_\sigma s_j$۔ که $t \in V(p)$ وي، نو وايو چې $p$
  په~$t$ کښې \emph{رښتيا دے}۔
\end{defn}

اړوند بدلون دا دے چې په مدل~$\mModel M$ کښې د وخت په ټکي~$t$
د !!a{formula} د رښتياوالي ارزونه د تاريخچې~$\sigma$ په نسبت کوو، چې
$t$ پکښې يو حالت وي۔ د !!{propositional
variable}s او د
رښتياوالي تابع نښلوونکو معنٰی پوهنه بدلولو ته اړتيا نشته۔ هغه ټول
د \olref[sem]{defn:tmodels} په شان عين دي، ځکه هېڅ يو يې
~$\sigma$ ته رجوع نۀ کوي۔ خو اوس د راتلونکي عامل~$\Ftemp$
معنا بيا تعريفوو او د دې تاريخچو په نسبت مودالي عامل~$\Diamond$
ورزياتوو۔

\begin{defn}\ollabel{defn:phmodels}
  په~$\mModel M$ کښې په~$t, \sigma$ باندې \emph{د !!a{formula}~$!A$
  رښتياوالے}، چې په نښو کښې ئې $\mSat{M}{!A}[t, \sigma]$ ليکو:
  \begin{enumerate}
  \item\ollabel{defn:sub:phmodels-f} \indcase{!A}{\Ftemp
    !B}{$\mSat{M}{\indfrm}[t, \sigma]$ هله او يوازې هله چې د کوم
    $t' \in T$ دپاره $t \prec_\sigma t'$ او
    $\mSat{M}{!B}[t', \sigma]$ وي۔}
  \item\ollabel{defn:sub:phmodels-diamond} \indcase{!A}{\Diamond
    !B}{$\mSat{M}{\indfrm}[t, \sigma]$ هله او يوازې هله چې د
    کومې $\sigma' \in C$ دپاره چې $t$ پکښې راځي،
    $\mSat{M}{!B}[t, \sigma']$ وي۔}
  \end{enumerate}
\end{defn}

نور زماني او مودالي عاملونه هم په ورته ډول تعريفولے شو۔ اوس بيا هغه
دعوې هم څرګندولے شو چې زمان او موداليت سره يوځاے کوي۔ د بېلګې په
توګه، «$p$ به واقع نۀ شي، خو په بله ممکنه تاريخچه کښې ښايي واقع شي»
په $\lnot \Ftemp p \land \Diamond \Ftemp p$ ښودلے شو۔ دا په هغه
ټکي او تاريخچه کښې رښتيا وي چې $p$ پکښې په راتلونکي حالت کښې رښتيا
نۀ شي، خو بله تاريخچه شته چې $p$ به پکښې رښتيا شي۔ د منجمدې سرچينې
د مثال په انګرېزي عبارت کښې «ښايي واقع شوے واے» تېر مهال راغلے؛
خو فورمول او ورپسې تشريح دواړه د بلې تاريخچې راتلونکے امکان ښيي۔

\end{document}
